Modern site reliability engineering teams operate mature anomaly detection pipelines; what is missing is a fast bridge from a fresh telemetry signal to the past incident that explains it. We build and characterize a retrieval-augmented triage system that, given a 5-minute window of logs, traces, and Kubernetes events, retrieves the most-similar past Jira tickets from a memory of engineer-voice incident reports.

Our analysis on a synthetic-but-realistic dataset of 2940 test windows and 347 prior tickets yields five empirical findings. \emph{First}, retrieval quality scales monotonically with deployment history: as the number of memory-resident prior tickets compatible with a window grows from 1--2 to 21+, Hit@1 climbs from 0.048 to 0.250 ($5.2\times$) and MRR climbs from 0.095 to 0.250 ($2.6\times$). Telemetry-only models (gradient-boosted classifier on 94 numeric features) cannot exhibit this curve --- they have no retrieval head. \emph{Second}, anomaly detection and retrieval are orthogonal: a numeric baseline already saturates the binary anomaly task at PR-AUC 0.77, fifteen percentage points ahead of our retrieval-augmented system, and we report this trade-off honestly. \emph{Third}, fine-tuning a domain-specific cross-encoder reranker on $(\textit{window}, \textit{ticket})$ pairs shifts probability mass into Hit@5 (+10--20\% relative) at the cost of top-1 precision, implying that ``review top-K candidates'' is a better product framing than ``auto-cite top-1.'' \emph{Fourth}, the system is robust to memory noise: with 50\% of memory occupied by plausible-but-fake distractor tickets, Hit@5 is exactly invariant across a four-point sweep and top-1 erodes only 4\% relative. \emph{Fifth}, cold-start novelty detection is unsolved in the current system --- of 333 cold-start anomalies, zero are flagged correctly --- a clear failure mode motivating future work.

Methodologically, we identify a subtle but important correction: the standard $\text{Recall@}K = |\text{top-}K \cap \text{gold}| / |\text{gold}|$ definition is mechanically capped at $K / |\text{gold}|$, making the depth-stratified curve appear non-monotone. We argue Hit@K is the right ``did the engineer find a relevant ticket'' metric. Translated to engineer-time-saved, a memory-augmented system with our fine-tuned reranker cuts mean diagnosis time from 30~minutes (telemetry-only baseline) to 23.6~minutes per resolvable incident.
